\documentclass[11pt]{article}
\usepackage{amsmath,amsthm,amssymb}
\usepackage[margin=1.1in]{geometry}
\usepackage{booktabs,longtable,array}
\usepackage[hidelinks]{hyperref}
\usepackage{microtype}

\newtheorem{theorem}{Theorem}[section]
\newtheorem{proposition}[theorem]{Proposition}
\newtheorem{lemma}[theorem]{Lemma}
\newtheorem{corollary}[theorem]{Corollary}
\theoremstyle{definition}
\newtheorem{definition}[theorem]{Definition}
\theoremstyle{remark}
\newtheorem{remark}[theorem]{Remark}

\newcommand{\Z}{\mathbb{Z}}
\newcommand{\Q}{\mathbb{Q}}
\newcommand{\N}{\mathbb{N}}
\newcommand{\pr}{\varrho}
\newcommand{\gam}{\Gamma}
\DeclareMathOperator{\Nm}{N}
\DeclareMathOperator{\sgn}{sign}
\newcommand{\mat}[4]{\left(\begin{smallmatrix}#1&#2\\#3&#4\end{smallmatrix}\right)}

\title{Odd/Odd Records for a Relative Approximation Error:\\
Three Quadratic Classifications and an Exact Certificate}
\author{Elias De Jes\'us\\
\small ORCID: \href{https://orcid.org/0009-0007-0190-9143}{0009-0007-0190-9143}}
\date{7 November 2026}

\begin{document}
\maketitle

\begin{abstract}
For a quadratic irrational $\alpha>1$ and positive \emph{odd} integers $m,k$, set
$H=m\alpha+k$ and $\varepsilon=|m\alpha-k|/(m\alpha+k)$. A state $(m,k)$ is a
\emph{strict chronological record} if $\varepsilon$ is smaller than at every admissible state
of smaller $H$. We classify the strict records completely for $\alpha=\varphi$,
$\alpha=1+\sqrt3$ and $\alpha=\sqrt3$. In each case they are a finite explicit union of
orbits of a single unimodular matrix, with signed defect cycles $(-1,+1,-5)$, $(-3,+1,-11)$
and $(-2)$ respectively. The proofs rest on two exact identities that make $\varepsilon$ a
function of the defect $N$ and the conjugate factor $\eta$ alone, on a norm cutoff
$|N|\le 18,204,27$ derived from heights, and on a finite exact certificate that decides every
remaining comparison. Two phenomena require care and are isolated: at $\varphi$ the
competition between the Fibonacci and Lucas ladders is a genuine tie at leading order, decided
by an integer identity $N_B=-5N_A$; and at $\sqrt3$ a companion family has \emph{exactly}
equal error at larger height, so it is excluded by equality rather than by loss. The results
resolve two conjectures of the author's earlier note and correct one assertion there. Nothing
here is formally verified or refereed, and the uniform generalisation is left open.
\end{abstract}

\section{Introduction}\label{sec:intro}

Let $\alpha>1$ be irrational and consider pairs $(m,k)$ of positive odd integers. Write
\begin{equation}\label{eq:Heps}
H=H(m,k)=m\alpha+k,\qquad
\varepsilon=\varepsilon(m,k)=\frac{|m\alpha-k|}{m\alpha+k}.
\end{equation}
Thus $\varepsilon$ is a \emph{relative} error, normalised by the same quantity $H$ that orders
the states. Because $\alpha$ is irrational, distinct integer pairs have distinct $H$, so the
states are linearly ordered by height with no ties.

\begin{definition}[Records]\label{def:records}
A state $(m,k)$ is a \emph{strict chronological record} if
$\varepsilon(m,k)<\varepsilon(m',k')$ for every admissible state $(m',k')$ with $H(m',k')<H(m,k)$.
It is a \emph{weak record} if $\varepsilon(m,k)\le\varepsilon(m',k')$ for every such state.
Every strict record is therefore a weak record, and a weak record that is \emph{not} strict
is one for which equality holds at some lower state: it ties the running minimum instead of
improving it.
\end{definition}

Throughout, ``records'' without qualification means strict records. Section~\ref{sec:classify}
states three complete classifications of them. We also exhibit, at $\alpha=\sqrt3$, an explicit
infinite family of weak records that are not strict (Theorem~\ref{thm:C}\,(ii)); we do
\emph{not} claim to classify all weak records, and in particular we do not treat imprimitive
weak records.

\paragraph{What this note does.}
The three classifications are proved, not merely computed to a large height. The architecture
is:
\begin{enumerate}\itemsep2pt
\item two exact identities (Lemma~\ref{lem:master}) reduce $\varepsilon$ and $H$ to functions
of the defect $N$ and the conjugate factor $\eta$;
\item imprimitive states are discarded (Lemma~\ref{lem:imprim}), the admissible unit step is
identified (Section~\ref{sec:lattice}), and a \emph{height-only} cutoff bounds $|N|$ on every
record (Theorem~\ref{thm:cutoff});
\item the finitely many surviving unit orbits are enumerated exactly inside a bounded lattice
rectangle, with no unique-factorisation hypothesis (Proposition~\ref{prop:enum});
\item every cross-orbit comparison of error and of height is reduced to the sign of one
algebraic quantity, constant along the orbits (Section~\ref{sec:compare}); and
\item a finite exact check covers the initial segment, beyond which the pattern is proved
to repeat (Section~\ref{sec:tail}).
\end{enumerate}
A machine-generated table of every orbit, its classification, and a checkable exclusion
witness is given in Appendix~\ref{app:tables}.

\paragraph{Relation to the author's earlier note.}
This is a companion to \cite{DeJesus2026}, pinned at
\begin{center}\small\ttfamily
commit 10ea1c62c6866990ded0a85a531883841d38d034\\
PDF SHA-256 f36c23ab099a7a0bca1e68dfe699ea02e15ef6e767bbc52428a77b272da247b3
\end{center}
\noindent In that note the
$\sqrt2$ analogue is proved, while the $\varphi$ and $1+\sqrt3$ chronologies are
\emph{computationally verified} with completeness left conjectural, and the note states
explicitly that its norm cutoff for those two slopes is assumed rather than derived.
Section~\ref{sec:earlier} records precisely which statements are resolved and which one is
corrected. (Citations there use the numbering of the pinned PDF: the two conjectures are
Conjecture 5.11 and Conjecture 5.13, and the competition remark is Remark 5.10. The
repository \texttt{README} refers to them as ``Conjectures 5.8, 5.10''; that is a
\texttt{README} error, not a different revision.)

\paragraph{What this note does not do.}
It does not prove a uniform theorem for general quadratic slopes and congruence classes;
Section~\ref{sec:outlook} explains where the present argument uses $t\ge0$, $\alpha'<0$ and the
odd/odd class specifically, and gives a counterexample to a naive generalisation. It does not
determine novelty relative to the literature beyond the limited comparisons in
Section~\ref{sec:related}. Nothing here is machine-checked in a proof assistant, nor refereed.

\section{Setting and the two master identities}\label{sec:master}

Let $\alpha$ satisfy $\alpha^2-t\alpha-s=0$ with $t,s\in\Z$, let $\alpha'$ be its conjugate,
and put
\[
D_0=\alpha-\alpha'>0,\qquad
\xi=k-m\alpha,\qquad \eta=k-m\alpha',\qquad
N=\xi\eta=k^2-tkm-sm^2 .
\]
$N$ is the integral defect. For the three slopes of this note:
\[
\begin{aligned}
\varphi:\quad &t=s=1, &&D_0=\sqrt5, &&N=k^2-km-m^2;\\
1+\sqrt3:\quad &t=s=2, &&D_0=2\sqrt3, &&N=k^2-2km-2m^2;\\
\sqrt3:\quad &t=0,\ s=3, &&D_0=2\sqrt3, &&N=k^2-3m^2 .
\end{aligned}
\]
In all three, $t\ge0$ and $\alpha'<0$, so $\eta>0$ on every positive state. Set
\begin{equation}\label{eq:G}
G=2\alpha\eta^2-tN .
\end{equation}

\begin{lemma}[Master identities]\label{lem:master}
For every state with $\eta>0$,
\begin{equation}\label{eq:masters}
H=\frac{G}{D_0\,\eta},\qquad
\varepsilon=\frac{|N|\,D_0}{G},\qquad\text{and in particular}\qquad
G>0 .
\end{equation}
Thus $\varepsilon$ and $H$ are determined by the pair $(N,\eta)$ alone.
\end{lemma}

\begin{proof}
From $\xi=k-m\alpha$ and $\eta=k-m\alpha'$ we get $\eta-\xi=mD_0$ and
$\alpha\eta-\alpha'\xi=kD_0$, hence $m=(\eta-\xi)/D_0$ and $k=(\alpha\eta-\alpha'\xi)/D_0$, so
\[
H=m\alpha+k=\frac{\alpha(\eta-\xi)+\alpha\eta-\alpha'\xi}{D_0}
=\frac{2\alpha\eta-(\alpha+\alpha')\xi}{D_0}=\frac{2\alpha\eta-t\xi}{D_0}.
\]
Substituting $\xi=N/\eta$ gives $H=(2\alpha\eta^2-tN)/(D_0\eta)=G/(D_0\eta)$, and $H>0$ forces
$G>0$. Then $\varepsilon=|\xi|/H=(|N|/\eta)\cdot D_0\eta/G=|N|D_0/G$.
\end{proof}

\begin{remark}[The error score, and why height is separate]\label{rem:score}
Since $\varepsilon=|N|D_0/G$, records are the running minima of $|N|/G$, equivalently the
running maxima of
\[
S:=\frac{G}{|N|}=2\alpha W-t\,\sgn(N),\qquad W:=\frac{\eta^2}{|N|} .
\]
So the \emph{error score} is $2\alpha W$ corrected by a bounded term depending only on the
\emph{sign} of $N$. The height order is a different question and is treated separately
(Section~\ref{sec:compare}): by \eqref{eq:masters}, $H=G/(D_0\eta)$ depends on $N$ as well as
$\eta$, so it is not a function of $S$, and in particular $\eta$-order is \emph{not}
height-order across different defects. Example: at $\alpha=\varphi$ the primitive odd/odd
states $(3,1)$ and $(1,3)$ have
$\eta=3\varphi-2\approx2.8541<3.618\approx\varphi+2=\eta$ respectively, while
$H=3\varphi+1\approx5.8541>4.618\approx\varphi+3$. Any argument that infers one order from the
other must be checked.
\end{remark}

\begin{lemma}[Imprimitive states never record]\label{lem:imprim}
If $g=\gcd(m,k)>1$ then $(m/g,k/g)$ is again an odd state, has the same $\varepsilon$, and has
strictly smaller $H$. Consequently no imprimitive state is a strict record, and $N$ may be
assumed primitively represented.
\end{lemma}

\begin{proof}
$m,k$ odd forces $g$ odd, so $m/g$ and $k/g$ are odd and the reduced pair is admissible. Both
$|m\alpha-k|$ and $m\alpha+k$ scale by $g$, so $\varepsilon$ is unchanged, while
$H$ is divided by $g>1$. A strict record must beat every state of smaller height, in
particular this one, which it cannot.
\end{proof}

The proof of Lemma~\ref{lem:imprim} uses that the gcd of two odd numbers is odd. This is
special to the odd/odd class; see Section~\ref{sec:outlook}.

\section{The admissible lattice and the admissible unit}\label{sec:lattice}

Write $\eta=a+b\omega$ with $\omega=\varphi$ or $\sqrt{D}$ as appropriate, so that the state is
recovered by
\[
(m,k)=(b,\,a+b)\ \text{ for }\varphi\text{ and }1+\sqrt3,\qquad (m,k)=(b,\,a)\ \text{ for }\sqrt3 .
\]
Admissibility ($m,k$ both odd) becomes a congruence on $(a,b)$ modulo $2$: ``$a$ even, $b$
odd'' in the first two cases, ``$a$ odd, $b$ odd'' for $\sqrt3$.

\begin{proposition}[Admissible unit step]\label{prop:unit}
Let $u$ be the fundamental unit of the relevant order and let $U\in GL_2(\Z)$ be the matrix of
multiplication by $u$ on the basis $(1,\omega)$. Reducing modulo $2$ places $U$ in the finite
group $GL_2(\Z/2)$ of order $6$, so some power $U^j$ fixes the admissibility class; the least
such $j$ gives the \emph{admissible unit} $\pr=u^j$. Multiplication by $\pr$ preserves
admissibility and primitivity, and multiplies $\eta$ by $\pr$ while multiplying $N$ by
$\delta:=\Nm(\pr)$. On the column $(m,k)$ it acts by a matrix with positive entries:
\[
\varphi:\ \pr=\varphi^3=1+2\varphi,\ \ \begin{pmatrix}1&2\\2&3\end{pmatrix},\ \delta=-1;\qquad
1+\sqrt3:\ \pr=(2+\sqrt3)^2,\ \ \begin{pmatrix}3&4\\8&11\end{pmatrix},\ \delta=+1;
\]
\[
\sqrt3:\ \pr=2+\sqrt3,\ \ \begin{pmatrix}2&1\\3&2\end{pmatrix},\ \delta=+1 .
\]
\end{proposition}

The minimal $j$ can be strictly smaller than the order of $U$ modulo $2$, because the class
need only be \emph{fixed}, not acted on trivially: at $\sqrt3$ the class (odd,\,odd) is fixed by
the swap $\mat{0}{1}{1}{0}$, so $j=1$ although $U\not\equiv I$. The values are $j=3,2,1$.

\begin{remark}[$\varphi^3$ versus $\varphi^6$]\label{rem:phi3}
The matrix $\mat{1}{2}{2}{3}$ sends $(1,1)\mapsto(3,5)$ and is the $\varphi^3$ step. The
matrix $\mat{5}{8}{8}{13}$, which appears in \cite{DeJesus2026}, sends $(1,1)\mapsto(13,21)$
and is the $\varphi^6$ step, i.e. the step \emph{within} a fixed-defect orbit. Both occur below
and must not be interchanged.
\end{remark}

\begin{lemma}[Signed transport]\label{lem:transport}
Along an orbit, with $\pr$ applied $p$ times to a chosen representative,
\begin{equation}\label{eq:transport}
N_p=\delta^p\,N_{\mathrm{rep}},\qquad \delta=\Nm(\pr).
\end{equation}
At $\varphi$, $\delta=-1$, so the \emph{sign} of $N$ alternates along a $\pr$-orbit and the
fixed-defect orbits are the $\pr^2=\varphi^6$ orbits. Taking $|N|$ therefore discards exactly
the parity of $p$ relative to the representative. Every comparison below that needs an
orientation transports the signed $N$ by \eqref{eq:transport}; magnitudes are used only where
the quantity is $\delta$-invariant.
\end{lemma}

Also useful later: the ladders have positive entries, so once a state has $m,k>0$ all its
successors do, and their heights increase without bound. Each orbit therefore has a finite
initial segment of non-positive states followed by an infinite positive tail.

\section{The height-only cutoff}\label{sec:cutoff}

Fix a \emph{comparison ladder}: an admissible orbit of known defect $\pm L$, whose positive
members we denote $A,C,\dots$ in increasing height. For $\varphi$ and $1+\sqrt3$ we take
$L=1$; for $\sqrt3$ no admissible state has $|N|=1$, because odd $m,k$ force
$k^2-3m^2\equiv1-3\equiv6\pmod 8$, and we take $L=2$ with the orbit of $(1,1)$.

\begin{theorem}[Cutoff]\label{thm:cutoff}
Let $B$ be a strict record whose height exceeds that of some positive ladder member, and let
$A,C$ be consecutive positive ladder members with $H_A<H_B\le H_C$. Put $Y=\eta_A^2$, $n=N_A$,
$L=|n|$. Then
\begin{equation}\label{eq:F}
|N_B|<\frac{H_C^2\,\varepsilon_A\,(D_0+t\varepsilon_A)}{2\alpha}=:F(Y,n),
\end{equation}
where
\[
\varepsilon_A=\frac{L\,D_0}{2\alpha Y-tn},\qquad
\frac{H_C}{H_A}=\frac{2\alpha \pr^2 Y-t\delta n}{\pr\,(2\alpha Y-tn)} .
\]
Consequently $|N_B|\le 18$ for $\alpha=\varphi$, $|N_B|\le204$ for $\alpha=1+\sqrt3$, and
$|N_B|\le27$ for $\alpha=\sqrt3$.
\end{theorem}

\begin{proof}
Write $z=(k-m\alpha)/H$, so $\varepsilon=|z|$. By Lemma~\ref{lem:master},
$\eta=H(D_0+tz)/(2\alpha)$ and hence
\begin{equation}\label{eq:Nheight}
|N|=|\xi|\eta=H\varepsilon\cdot\frac{H(D_0+tz)}{2\alpha}=\frac{H^2\varepsilon(D_0+tz)}{2\alpha}.
\end{equation}
The ladder heights increase to infinity, so consecutive members bracket $B$ as stated. Since
$B$ is a record, $\varepsilon_B<\varepsilon_A$; also $t\ge0$ and $|z_B|=\varepsilon_B$, so
$D_0+tz_B\le D_0+t\varepsilon_B$, and $x\mapsto x(D_0+tx)$ is strictly increasing for $x\ge0$.
Applying \eqref{eq:Nheight} to $B$ and then $H_B\le H_C$ gives \eqref{eq:F}. The displayed
formulas for $\varepsilon_A$ and $H_C/H_A$ follow from \eqref{eq:masters} with $\eta_C=\pr\eta_A$,
$N_C=\delta n$. Only heights were compared.

It remains to bound $F$ over the \emph{whole} ladder. Put
\[
u=\frac{t}{2\alpha Y}\quad(\text{and }u=\tfrac1{2\alpha Y}\text{ at }\varphi,\text{ where }L=t=1).
\]
Two cases. If $t=0$ then $u=0$ identically, $\varepsilon_A=LD_0/(2\alpha Y)$ and
$H_C/H_A=\pr$, so \eqref{eq:F} collapses to $F=L\pr^2$, a constant: no bound on $u$ is needed
or available, and this is the only case used at $\sqrt3$. If $t>0$ then $0<u<1$: at the first
ladder member $2\alpha Y>t$ (an exact inequality, checked in the certificate), and $Y$ is
multiplied by $\pr^2>1$ at each step while $t,\alpha$ are fixed, so $u$ strictly decreases.
The three closed forms with a $u$ in them are used only on $0<u<1$.
Substituting into \eqref{eq:F} and simplifying:
\[
\varphi,\ n=-1:\quad F=\pr^2\Bigl(1-\frac{u}{\pr^2}\Bigr)^{\!2}\frac{1+2u}{(1+u)^2}<\pr^2,
\]
because $(1-u/\pr^2)^2<1$ and $(1+2u)<(1+u)^2$ for $u>0$; and $\pr^2=\varphi^6=17.9442719\ldots$
\[
\varphi,\ n=+1:\quad F=\pr^2\frac{(1+u/\pr^2)^2}{(1-u)^2},
\]
which is increasing in $u$, hence decreasing in $Y$; its maximum over the ladder is attained at
the first member with $n=+1$, namely $Y=\varphi^8$, where $F=18.1960225\ldots<19$.
\[
1+\sqrt3:\quad F=\pr^2\Bigl(\frac{1-u/\pr^2}{1-u}\Bigr)^{\!2},
\]
increasing in $u$, hence maximal at the first ladder member $Y=(2+\sqrt3)^2$, where
$F=204.5520518\ldots<205$.
\[
\sqrt3:\quad t=0\ \Longrightarrow\ F=L\pr^2=2(2+\sqrt3)^2=27.8564064\ldots<28 .
\]
Taking integer parts gives the three bounds. The endpoint inequalities are decided exactly in
the certificate; the decimals are display only.
\end{proof}

\begin{remark}[Why the signs must be separated]\label{rem:signsep}
At $\varphi$ the two branches move in \emph{opposite} directions along the ladder: the $n=-1$
values increase ($17.5116\ldots,17.9435\ldots,17.9443\ldots$) while the $n=+1$ values decrease
($18.1960\ldots,17.9450\ldots,17.9443\ldots$), both towards $\pr^2$. So $F$ is not monotone in
$Y$, and no finite sample of ladder states bounds the infinite ladder. The separation above
does: one branch is bounded uniformly by $\pr^2$, the other attains its maximum at an explicit
first member.
\end{remark}

Combining Theorem~\ref{thm:cutoff} with Lemma~\ref{lem:imprim} and the congruence restrictions
on $N$ leaves very little at $\varphi$. There $N=k^2-km-m^2$ is \emph{odd}, since $m,k$ are;
$4N=(2k-m)^2-5m^2$ forces $4N$ to be a square modulo $5$, i.e.\ $N\equiv0,1,4\pmod5$, which
among the odd $|N|\le18$ leaves $|N|\in\{1,5,9,11,15\}$; and $3$ is inert in $\Z[\varphi]$
(the discriminant $5$ is not a square modulo $3$), so $3\mid N$ forces $3\mid\gcd(m,k)$ and
kills $|N|=9,15$ by Lemma~\ref{lem:imprim}. The admissible defects left at $\varphi$ are
therefore $|N|\in\{1,5,11\}$, and the four orbits of Appendix~\ref{app:tables} realise exactly
these.

\section{Exhaustive enumeration of the surviving orbits}\label{sec:enum}

\begin{proposition}[Bounded enumeration, no unique factorisation]\label{prop:enum}
Let $|N|\le B$. Reduce $\eta$ by powers of $\pr$ until $1\le\eta<\pr$. Then
$|\xi|=|N|/\eta\le B$ automatically, so $(\eta,\xi)$ lies in the rectangle
$[1,\pr)\times[-B,B]$ in the two real embeddings, and since $m=(\eta-\xi)/D_0$ and
$k=(\alpha\eta-\alpha'\xi)/D_0$ the integer coordinates of the representative are explicitly
bounded. An exhaustive search over that box finds every orbit.
\end{proposition}

No class-group or principal-ideal hypothesis is used: this is the classical reduction theory of
the generalised Pell equation, see \cite{Conrad}, Theorems 1.1 and 2.1. Reduced representatives may
have $m\le0$ or $k\le0$; they must be retained, since the positive states appear later in the
same orbit. The resulting counts are
\[
\varphi:\ 4\ \text{orbits at }|N|\le18,\qquad
1+\sqrt3:\ 64\ \text{at }|N|\le204,\qquad
\sqrt3:\ 6\ \text{at }|N|\le27 .
\]

\section{Exact comparisons: error and height}\label{sec:compare}

Let $A,B$ be orbits with reduced representatives $r_A,r_B\in[1,\pr)$ and representative defects
$n_A,n_B$. Consider $A$ at period $p-d$ and $B$ at period $p$, with offset $d\ge0$.

\begin{proposition}[Error comparison]\label{prop:err}
With $\gam:=|N_A|\eta_B^2-|N_B|\eta_A^2$,
\begin{equation}\label{eq:crit}
\varepsilon_B<\varepsilon_A
\iff
2\alpha\gam>t\bigl(|N_A|N_B-|N_B|N_A\bigr),
\end{equation}
and $\gam=\pr^{2p-2d}\gam_d$ with
$\gam_d=|N_A|r_B^2\pr^{2d}-|N_B|r_A^2$ independent of $p$. Hence
$\sgn\gam=\sgn\gam_d$ is one exact sign per pair $(A,B)$ and offset $d$. The right-hand side of
\eqref{eq:crit} is bounded in absolute value by $2t|N_A||N_B|$, so if $\gam_d\neq0$ then
\eqref{eq:crit} is decided by $\sgn\gam_d$ as soon as
\begin{equation}\label{eq:thr}
2\alpha\,\pr^{2p-2d}\,|\gam_d|\;>\;2t\,|N_A|\,|N_B| ,
\end{equation}
a strict inequality which first holds at an explicitly computable period $P(A,B,d)$. The
threshold starts at $p=d$, where $\pr^{2p-2d}=1$.
\end{proposition}

\begin{proposition}[Degenerate case $\gam_d=0$]\label{prop:gamma0}
If $\gam_d=0$ then \eqref{eq:crit} reduces to the integer condition
$0>t(|N_A|N_B-|N_B|N_A)$ with the defects transported by \eqref{eq:transport}:
$N_A=\delta^{p-d}n_A$, $N_B=\delta^{p}n_B$. Hence
\begin{equation}\label{eq:tiesign}
\sgn(\varepsilon_B-\varepsilon_A)=\sgn\bigl(t(|N_A|N_B-|N_B|N_A)\bigr)\in\{-1,0,+1\},
\end{equation}
a three-way outcome. If $\delta=-1$ this depends on the parity of $p$, so such a pair can give
\emph{alternating} verdicts; if $t=0$ the right-hand side vanishes identically and the two
states have \emph{exactly equal} error at every period.
\end{proposition}

\begin{proposition}[Height comparison]\label{prop:height}
With $A$ at period $q$ and $B$ at period $q+d$, put
\[
U=2\alpha(r_B\pr^{d}-r_A)\,r_A\,r_B\pr^{d},\qquad
V=\delta^{q}\,t\bigl(\delta^{d}n_B r_A-n_A r_B\pr^{d}\bigr).
\]
Then $\sgn(H_B-H_A)=\sgn(\pr^{2q}U-V)$. Here $U$ is independent of $q$, and $V$ depends on $q$
only through $\delta^q$, hence only through the parity of $q$ when $\delta=-1$. If $U\neq0$ the
sign is decided by $\sgn U$ once $\pr^{2q}|U|>|V|$, at an explicitly computable period.
\end{proposition}

\begin{proof}
By \eqref{eq:masters}, $D_0H=2\alpha\eta-tN/\eta$, so
$D_0(H_B-H_A)=2\alpha(\eta_B-\eta_A)-t(N_B/\eta_B-N_A/\eta_A)$. Substitute
$\eta_A=r_A\pr^{q}$, $\eta_B=r_B\pr^{q+d}$ and the transported defects, then clear the positive
factor $\pr^{-q}r_Ar_B\pr^{d}$.
\end{proof}

\begin{proposition}[$U=0$ is an index shift, not a tie]\label{prop:U0}
$U=0$ if and only if $r_A=r_B\pr^{d}$. With distinct representatives in $[1,\pr)$ this is
impossible: for $d\ge1$ it forces $r_A\ge\pr$, and for $d=0$ it forces $r_A=r_B$, i.e. the same
orbit. Moreover $U=0$ would imply $V=0$: taking norms of $r_A=r_B\pr^d$ gives
$n_A=\delta^dn_B$, and substituting both identities annihilates $V$. Hence $U=0$ denotes the
same arithmetic state after an index shift, in contrast to $\gam_d=0$, which does occur between
distinct states.
\end{proposition}

When the two thresholds are combined one must convert period coordinates:
Proposition~\ref{prop:err} places $A$ at $p-d$ and $B$ at $p$, while
Proposition~\ref{prop:height} places $A$ at $q$ and $B$ at $q+d$, so $q=p-d$ and the height
threshold must be shifted by $+d$ before taking a maximum.

\section{From a finite check to the infinite tail}\label{sec:tail}

\begin{lemma}[Error monotonicity]\label{lem:mono}
Along any orbit, $\varepsilon$ strictly decreases. Indeed $\varepsilon=|N|D_0/G$ and one unit
step gives
\begin{equation}\label{eq:Ginc}
G_{\mathrm{next}}-G=2\alpha\eta^2(\pr^2-1)-tN(\delta-1).
\end{equation}
If $\delta=+1$ the second term vanishes and the first is positive, so $G$ increases at every
step unconditionally. If $\delta=-1$ the second term is bounded by $2t|N|$ while the first grows
like $\pr^{2p}$, so once $2\alpha\eta^2(\pr^2-1)>2t|N|$ it holds at every later step; the
finitely many earlier steps are checked directly.
\end{lemma}

\begin{proposition}[Only offsets $d\in\{0,1\}$ matter, in the tail]\label{prop:window}
Let $p_0$ be any period such that \emph{(a)} every orbit has a positive state at some period
$\le p_0-1$, and \emph{(b)} every pairwise ordering threshold of
Proposition~\ref{prop:height} is already decisive. Then for every orbit $B$ and every
$p\ge p_0$, and for each orbit $A$, the latest $A$-state strictly below the period-$p$ state of
$B$ in height occurs at period $p$ or $p-1$. If $B$ beats both of those, it beats every earlier
state of $A$.
\end{proposition}

\begin{proof}
Three ingredients. (i) Normalisation: all representatives lie in $[1,\pr)$, so for any orbits
$A,A'$ one has $r_A\pr^{p-1}<r_{A'}\pr^{p}$, because $r_A/r_{A'}<\pr$; hence the period-$(p-1)$
state of any orbit precedes the period-$p$ state of any other in the variable $\eta$.
(ii) Certified height order: the \emph{height} comparison is supplied by
Proposition~\ref{prop:height}, pairwise and exactly --- not inferred from $\eta$
(Remark~\ref{rem:score}) --- and hypothesis (b) is exactly what makes its sign the constant
$\sgn U$. (iii) Monotonicity: by Lemma~\ref{lem:mono}, a state of orbit $A$ at period
$q\le p-2$ has larger $\varepsilon$ than the period-$(p-1)$ state of $A$, which by (i) and (ii)
lies below $B$ and which $B$ already beats. Hypothesis (a) guarantees that the period-$(p-1)$
state of $A$ exists and is positive, so the window is populated.
\end{proof}

\begin{remark}[Both hypotheses are needed]\label{rem:windowfails}
Without (a) the conclusion is false, and not only vacuously. At $\alpha=1+\sqrt3$ take $B$ to
be the orbit through $(7,1)$, of defect $-111$, at its period $p=0$, where
$H=8+7\sqrt3=20.1244\ldots$. The orbit of defect $+193$ has no positive state at period $0$;
its first is $(1,15)$ at period $1$, with $H=16+\sqrt3=17.7320\ldots<H(7,1)$. So the latest
competitor state below $B$ sits at period $p+1$, outside the window. There are $13$ such pairs
at $1+\sqrt3$, all at $p=0$, and none at $\varphi$ or $\sqrt3$. Every one of them lies inside
the region $H\le302$ that the unpruned enumeration of Section~\ref{sec:tail} covers, which is
why the global tail start is raised to the least $p_0$ satisfying (a) and (b) --- the values
$p_0=2,2,3$ in the table below.
\end{remark}

Consequently the classification of an orbit is determined by finitely many exact signs: for each
competitor and each $d\in\{0,1\}$, the sign of $\gam_d$, or the three-way rule
\eqref{eq:tiesign} when $\gam_d=0$, together with the height order of
Proposition~\ref{prop:height}. Each orbit receives one of four verdicts:
\emph{strict} (beats every earlier competitor at every period parity), \emph{alternating}
(beats them at one parity, loses at the other), \emph{equality} (ties exactly with some earlier
competitor, hence not a strict record), or \emph{excluded}.

\paragraph{The finite segment.}
Everything below the period at which all those signs become decisive is covered by
\emph{unpruned} enumeration: all odd pairs with $H$ below an exact bound, with no defect
restriction. Take for $p_0$ the least period satisfying both hypotheses of
Proposition~\ref{prop:window}: the maximum over all pairs and offsets of the error and
ordering thresholds above (hypothesis (b)), raised until every orbit has a positive state at
some period $\le p_0-1$ (hypothesis (a)). This gives:
\begin{center}
\begin{tabular}{lccc}
\toprule
$\alpha$ & tail start $p_0$ & largest pre-tail $H$ & unpruned cap \\
\midrule
$\varphi$ & $2$ & $(29+7\sqrt5)/2=22.3262\ldots$ & $23$ \\
$1+\sqrt3$ & $2$ & $206+55\sqrt3=301.2628\ldots$ & $302$ \\
$\sqrt3$ & $3$ & $23+13\sqrt3=45.5167\ldots$ & $46$ \\
\bottomrule
\end{tabular}
\end{center}
\noindent The table is generated by \texttt{certificate/tail\_table.py}, which also reports the
state at which each maximum is attained. The unpruned enumerations below these caps (41, 4173 and 152 admissible states respectively)
agree with the claimed families, and the surviving orbits are identified by their actual
recurrences, not merely by their defect labels. Larger unpruned brute-force runs were also performed --- to
$H\le3000$ at $\varphi$ and $1+\sqrt3$ and to $H\le2000$ at $\sqrt3$, each cross-checked against
an independently pruned scan, and the record families compared with their recurrences out to
$H\le10^{12}$ --- but these are supporting verification, not part of the argument.

\section{The three classifications}\label{sec:classify}

Let $F_0=0,F_1=1$ and $L_0=2,L_1=1$.

\begin{theorem}[$\alpha=\varphi$]\label{thm:A}
The strict chronological records are exactly the states
\[
(F_{3j+1},F_{3j+2}),\quad j\ge0,\qquad\text{together with}\qquad (L_{6j+4},L_{6j+5}),\quad j\ge0,
\]
merged in increasing $H$. The signed defect cycle along the merged sequence is
$(-1,+1,-5)$, repeating. The first ladder is the orbit of $(1,1)$ under $\varphi^3$,
i.e. under $\mat{1}{2}{2}{3}$; within it the defect alternates between $-1$ and $+1$, and the
fixed-defect step is $\varphi^6=\mat{5}{8}{8}{13}$.
\end{theorem}

\begin{theorem}[$\alpha=1+\sqrt3$]\label{thm:B}
The strict chronological records are exactly the three orbits of
$\mat{3}{4}{8}{11}$ seeded at $(1,1)$, $(1,3)$ and $(5,13)$, merged in increasing $H$.
Their defects are constant along each orbit, equal to $-3$, $+1$ and $-11$, and the merged
defect cycle is $(-3,+1,-11)$, repeating. Here $\pr=7+4\sqrt3$ and $\delta=+1$.
\end{theorem}

\begin{theorem}[$\alpha=\sqrt3$]\label{thm:C}
\emph{(i)} The strict chronological records are exactly the orbit of $(1,1)$ under
$\mat{2}{1}{3}{2}$, that is
\[
(1,1),(3,5),(11,19),(41,71),(153,265),(571,989),\dots,
\]
all of defect $-2$. Here $t=0$, so $H=\eta$ and $h^2\varepsilon=\tfrac12$ exactly on every
strict record, where $h=H/2$.
\emph{(ii)} For each such state $A=(m,k)$ the companion $B=(k,3m)$ satisfies
\[
\eta_B=\sqrt3\,\eta_A,\qquad N_B=+6,\qquad H_B=\sqrt3\,H_A,\qquad \varepsilon_B=\varepsilon_A ,
\]
and $B$ is again an admissible primitive state. Moreover $B$ lies strictly between $A$ and the
next strict record $\pr A$:
\begin{equation}\label{eq:companionwindow}
H_A<\sqrt3\,H_A<(2+\sqrt3)H_A=H_{\pr A},
\end{equation}
so the running minimum of $\varepsilon$ over all states of height $<H_B$ is $\varepsilon_A$
itself. Hence the companion family $(1,3)$, $(5,9)$, $(19,33)$, $(71,123)$,
$(265,459),\dots$ consists of \emph{weak} records that are not strict.
\end{theorem}

\begin{proof}[Proof of (ii)]
$\eta_A=k+m\sqrt3$ and $\sqrt3\,\eta_A=3m+k\sqrt3=\eta_B$ with $(m_B,k_B)=(k,3m)$. Norms
multiply: $N_B=\Nm(\sqrt3)N_A=(-3)(-2)=6$. Since $t=0$, $H=\eta$ by \eqref{eq:masters}, so
$H_B=\sqrt3H_A$; and $\varepsilon=|N|/H^2$ gives
$\varepsilon_B=6/(3H_A^2)=2/H_A^2=\varepsilon_A$. Admissibility: $k$ is odd and $3m$ is odd.
Primitivity: $k^2-3m^2=-2$ forces $3\nmid k$ (else $0\equiv1\bmod3$) and $\gcd(k,m)=1$, so
$\gcd(k,3m)=1$. Finally \eqref{eq:companionwindow} holds because $1<\sqrt3<2+\sqrt3$, and by part~(i) the
strict records are exactly the $\pr$-orbit of $A$, so no state of height below $H_B$ has error
smaller than $\varepsilon_A$. Since $\varepsilon_B=\varepsilon_A$ and $H_B>H_A$, the state $B$
meets the running minimum without improving it: it is a weak record and not a strict one.
\end{proof}

We emphasise that Theorem~\ref{thm:C}\,(ii) exhibits one infinite family of non-strict weak
records; it is \emph{not} a classification of all weak records, and imprimitive weak records are
not considered.

\subsection{\texorpdfstring{The $\varphi$ tie, resolved by an integer}{The phi tie, resolved by an integer}}\label{sec:tie}

The $|N|=1$ admissible states are exactly $(F_n,F_{n+1})$ with $n\equiv1\pmod3$: $\eta$ is then
a unit $\pm\varphi^{\,n}$, the class condition modulo $2$ forces $n\equiv1\pmod3$, and $F_n$ and
$F_{n+1}$ are both odd precisely when $3\nmid n(n+1)$. The $|N|=5$ admissible states are
$\eta=\pm\sqrt5\,\varphi^{\,n}$ with the same congruence on $n$, since $5$ ramifies and
$(\sqrt5)=(2\varphi-1)$ is the only ideal of norm $5$; their defect is $N=-5(-1)^n$.

For such $\eta$ we abbreviate $\varepsilon(\eta)$ for the error of the corresponding state, by
Lemma~\ref{lem:master} a function of $(N,\eta)$.

\begin{theorem}[The tie and its resolution]\label{thm:tie}
For $n\equiv1\pmod3$ put $A$ the unit state at $\eta_A=\varphi^{\,n}$ and $B$ the norm-$5$ state
at $\eta_B=\sqrt5\,\varphi^{\,n}$. Then $\eta_B/\eta_A=\sqrt5$, so
$(\eta_B/\eta_A)^2=5=|N_B|$: the leading terms cancel exactly and
\begin{equation}\label{eq:tiepair}
\varepsilon(\varphi^{\,n})=\frac{\sqrt5}{2\varphi^{\,2n+1}-(-1)^n},\qquad
\varepsilon(\sqrt5\,\varphi^{\,n})=\frac{\sqrt5}{2\varphi^{\,2n+1}+(-1)^n}.
\end{equation}
Moreover $N_A=(-1)^n$ and $N_B=-5N_A$, so by \eqref{eq:tiesign} $B$ beats $A$ if and only if
$N_A=+1$, i.e. $n$ even. Among $n\equiv1\pmod3$ this is $n\equiv4\pmod6$, and those states are
exactly the Lucas pairs $(L_{6j+4},L_{6j+5})$.
\end{theorem}

So the competition is decided not by a multiplicative gap but by the bounded term $\pm1$ in
\eqref{eq:tiepair}, equivalently by the integer identity $N_B=-5N_A$. The normalised
representative of the norm-$5$ orbit has $N_{\mathrm{rep}}=+5$; the states that actually win are
those with $N=-5$, reached at the other period parity. This is the ``asymptotically degenerate
competition'' of \cite[Remark 5.10]{DeJesus2026}.

\begin{remark}[Method note]\label{rem:signmethod}
Two of the comparisons above are invisible to $|N|$: the alternating verdict of
Proposition~\ref{prop:gamma0} when $\delta=-1$, and the $V$-term of
Proposition~\ref{prop:height}. In the course of this work, asking of each computed quantity
which distinctions it discards --- and whether a later step needs them --- located exactly these
two places, and the certificate now transports signed defects by \eqref{eq:transport} and uses
magnitudes only where the quantity is $\delta$-invariant. The mechanism itself is arithmetic:
$N_B=-5N_A$, and the exact cancellation in \eqref{eq:tiepair}.
\end{remark}

\subsection{Classification summary}\label{sec:summary}

\begin{center}
\begin{tabular}{lrrcccc}
\toprule
$\alpha$ & cutoff & orbits & strict & alternating & equality & excluded \\
\midrule
$\varphi$     & $18$  & $4$  & $1$ & $1$ & $0$ & $2$ \\
$1+\sqrt3$    & $204$ & $64$ & $3$ & $0$ & $0$ & $61$ \\
$\sqrt3$      & $27$  & $6$  & $1$ & $0$ & $1$ & $4$ \\
\bottomrule
\end{tabular}
\end{center}

At $\varphi$ the alternating orbit is the norm-$5$ (Lucas) family of Theorem~\ref{thm:tie}; at
$\sqrt3$ the equality orbit is the companion family of Theorem~\ref{thm:C}\,(ii). Appendix
\ref{app:tables} lists every orbit with its representative, signed defect, verdict, and --- for
each excluded orbit --- a competitor, an offset, the deciding sign and the period from which it
is certified.

\section{Relation to the earlier note}\label{sec:earlier}

All references are to the pinned version of \cite{DeJesus2026} identified in
Section~\ref{sec:intro}.

\paragraph{Resolved.}
Conjecture 5.11 ($\varphi$-completeness) is Theorem~\ref{thm:A}; Conjecture 5.13
($(1+\sqrt3)$-completeness) is Theorem~\ref{thm:B}. The earlier note states that for these two
slopes its pruned search range is labelled computational ``precisely because the analogous bound
is'' those conjectures: that is, the norm cutoff was assumed. Theorem~\ref{thm:cutoff} derives
it.

Remark 5.10 of the baseline already identifies the mechanism of the degenerate competition
correctly --- that $|N|=5$ is compensated to leading order, that the decision falls to the
bounded term in $L_n^2-5F_n^2=4(-1)^n$, and that its sign selects the odd Lucas pairs with
$n\equiv4\pmod6$ --- and says explicitly that this is ``an observed pattern with an identified
finite mechanism, not a proved theorem''. Theorem~\ref{thm:tie} proves it, for every admissible
$n$ at once, from the closed forms \eqref{eq:tiepair}; and it places the phenomenon in a general
frame, as the case $\gam_d=0$ of Proposition~\ref{prop:gamma0}, where the same three-way integer
rule also produces the permanent equality at $\sqrt3$.

\paragraph{Corrected.}
Observation 5.9 of the baseline describes the $\varphi$ chronology as having ``three records
per multiplicative period $\varphi^6$ in height'', and Observation 5.12 says the same of
$1+\sqrt3$ with period $(2+\sqrt3)^2$; both are labelled computationally verified. Exact
multiplicativity holds for $\eta$, not for $H$: by \eqref{eq:masters},
\begin{equation}\label{eq:heightper}
H_{\mathrm{next}}-\pr H=\frac{tN\,(\pr-\delta/\pr)}{D_0\,\eta},
\end{equation}
which vanishes only when $tN=0$. Concretely, with $\mat{3}{4}{8}{11}(1,1)=(7,19)$ and
$\varphi^6\cdot(1,1)=(13,21)$,
\[
H(7,19)-\pr H(1,1)=-8\sqrt3,\qquad
H(13,21)-\varphi^6H(1,1)=4-4\sqrt5,
\]
both nonzero. What is true is that $H_{\mathrm{next}}/H\to\pr$ along each orbit, by
\eqref{eq:heightper}, since $\eta\to\infty$; the baseline's observed record lists are not in
question, only the exactness of the height period. We do \emph{not} assert that every sufficiently large multiplicative height
window contains exactly three records; that is a statement about the interleaving of three
orbits with different multiplicative phases and would require a separate argument. At
$\alpha=\sqrt3$, where $t=0$, one has $H=\eta$ and the height period $2+\sqrt3$ \emph{is} exact
(Theorem~\ref{thm:C}).

\paragraph{New.}
The slope $\sqrt3$ does not appear in \cite{DeJesus2026}. It is the cleanest of the three: $t=0$
collapses the master identities to $H=\eta$ and $\varepsilon=|N|/H^2$, so $h^2\varepsilon$ is
exactly $|N|/4$, and it is also the case in which no admissible state has $|N|=1$, which is why
the cutoff is stated for a ladder of arbitrary known defect.

\section{The certificate}\label{sec:cert}

The finite content of the proofs is isolated in a small standalone bundle
: the directory \texttt{certificate/} of the companion
folder in the repository, documented by its own \texttt{README}. It is released under the tag
\begin{center}\texttt{congruence-records-v1.0.0}\end{center}
\noindent whose commit and per-file hashes are recorded in
\texttt{docs/RELEASE\_MANIFEST.md}. Two programs are run:
\texttt{certificate.py} prints the cutoff, the orbit enumeration, the classification and the
threshold data; \texttt{verify\_records.py} runs 31 checks and reports \texttt{VERIFY PASSED}.
The bundle uses the Python standard library only, reads no data files at runtime, and
recomputes everything from the definitions.

All ordering and comparison decisions are exact: arithmetic is carried out in $\Z[\sqrt D]$ with
integer or rational coordinates, and signs are decided by integer comparisons. Floating point
appears only as a presort key --- after which the sorted order is re-verified exactly --- and for
display. The checks include the master identities, the congruence restrictions on $N$ for each
slope, the imprimitivity reduction, the cutoff bound against every record found by unpruned
search, an independent recount of the orbits by direct rectangle enumeration, the four-way
classification against the claimed families, set equality of the record list with the
recurrences of Theorems~\ref{thm:A}--\ref{thm:C}, the signed-transport ledger
\eqref{eq:transport} and its invariance under change of representative, and the threshold data
of Section~\ref{sec:tail}.

\paragraph{Status of the verification.}
The certificate is exact arithmetic, not a formal proof: it is not machine-checked in a proof
assistant, and this note has not been refereed. Two independent audits of earlier versions of
the bundle were carried out with assistance of an automated reviewer; they re-derived the
classifications by separate orbit enumeration and tail logic, and they found and prompted the
repair of several defects in the written proof and in the implementation, including the
incorrect monotonicity justification repaired in Remark~\ref{rem:signsep} and a threshold
initialisation error. Those audits re-used this bundle's arithmetic module, so they are
independent in logic and in control flow but not in the underlying arithmetic implementation.
The release history in \texttt{docs/CHANGELOG.md} records the versions and what changed.

\section{Related work}\label{sec:related}

The objective studied here --- a relative error $|m\alpha-k|/(m\alpha+k)$, with the states
ordered by the same denominator $m\alpha+k$ --- is not one of the standard best-approximation
notions, and we make only comparisons that the cited definitions support.

Kim, Lee and Liao \cite{KLL2022} introduce the odd--odd continued fraction and prove that the
best \emph{$1$-rational} approximations of an irrational $x$ --- their Definition 1.1: $p/q$
with $p,q$ odd and coprime such that $|qx-p|<|bx-a|$ for every $a/b$ of the same parity type
with $0<b\le q$ --- are exactly its principal convergents (their Theorem 1.2). That objective is
$|qx-p|$ ordered by $q$, which is not \eqref{eq:Heps}; the two notions select different sets.
For instance at $\alpha=\varphi$ the state $(7,11)$ is a strict record here, while
$|3\varphi-5|=0.1459\ldots<0.3262\ldots=|7\varphi-11|$ with $3\le7$, so $11/7$ is not a best
$1$-rational approximation. A companion paper \cite{KLL2025} studies the three parity types and
best signed approximations with the same $|qx-p|$ objective. We make no claim that our
classifications do or do not follow from that machinery; establishing either would require an
argument we have not given.

Hurwitz-type theorems \emph{with} congruence conditions on numerator and denominator are a
separate line. Elsner's introduction \cite{Elsner} records it as beginning with Hartman, who
first imposed congruence conditions on $u$ and $v$, with a weaker theorem of Koksma (1951) and
with Uchiyama's result the sharpest before his own. Elsner's Theorem 1.1 requires the modulus
$s$ to have an odd prime divisor, and he records that the corresponding Corollary 1.1 fails for
$s=2$ with $h=1$ --- that is, for the odd/odd class --- giving $\xi=12+\sqrt{145}$ as a
witness. So the restriction studied here is the acknowledged
exceptional case of that line of work. Those results bound an approximation constant; they do
not enumerate a record set.

Cha and Kim \cite{ChaKim} determine the small values and the smallest accumulation point of the
Lagrange spectrum of intrinsic approximation on the unit circle, with height the denominator of
a primitive integer vector, and classify the very badly approximable points by Romik digit
expansions. That is an asymptotic spectrum on a different space with a different height, and we
draw no reduction in either direction. By contrast the accumulation values of
$h^2\varepsilon$ along the records here follow immediately from \eqref{eq:masters}:
$h^2\varepsilon\to\alpha|N|/(2D_0)$ along a fixed-defect orbit, which at $\sqrt3$ is exactly
$|N|/4$. These are values \emph{at} records and are not the same quantity as the staircase
envelope of $H\mapsto H^2\varepsilon$ between records.

The reduction theory used in Section~\ref{sec:enum} --- finiteness of the orbits of a given norm
under the unit group, and the bounded search that exhibits representatives --- is classical; see
\cite{Conrad}, Theorems 1.1 and 2.1. Our contribution is the three exact chronologies for the error and height
\eqref{eq:Heps}, the two equality mechanisms isolated in Theorems~\ref{thm:tie} and
\ref{thm:C}\,(ii), and the effective certificate that makes the passage from a finite
computation to the infinite tail explicit.

\section{Outlook, with its hypotheses stated}\label{sec:outlook}

It is natural to ask for a uniform statement: for a quadratic $\alpha$ and a congruence class on
$(m,k)$, are the strict records always a finite explicit union of unit orbits beyond an
effectively computable height? The present argument does not prove this, and three of its steps
use the particular setting.

\begin{itemize}\itemsep3pt
\item \textbf{Signs.} Theorem~\ref{thm:cutoff} uses $t\ge0$ and $\alpha'<0$, so that $\eta>0$ on
positive states and $x\mapsto x(D_0+tx)$ is increasing. Other sign patterns need separate
treatment.
\item \textbf{Imprimitivity.} Lemma~\ref{lem:imprim} uses that the gcd of two odd numbers is
odd, so dividing it preserves the class. This fails for a general class: at $\alpha=\sqrt6$ with
the class $(m,k)\equiv(1,1)\pmod3$, the state $(4,10)$ is a strict record within the class
although $\gcd=2$, and its reduction $(2,5)$ leaves the class. A general statement must quantify
over primitive states or retain imprimitive bounded-defect orbits.
\item \textbf{The comparison ladder.} Theorem~\ref{thm:cutoff} needs an admissible state of
\emph{known} defect $L$. Existence of some primitive seed in an ordinary residue pair
$(a,b)\bmod M$ with $\gcd(a,b,M)=1$ follows from a short construction: choose $m\equiv a$, let
$P$ be the product of the primes dividing $m$ but not $M$, and use the Chinese remainder theorem
to choose $k\equiv b\pmod M$ with $k\equiv1\pmod P$; common primes of $m$ and $M$ cannot divide
$b$, so $\gcd(m,k)=1$, and a large enough $k$ makes $\eta>0$. Its defect is nonzero and
computable. Obtaining a \emph{prescribed} or \emph{minimal} $|L|$ is a different and harder
question; the values $L=1,1,2$ used here were found by inspection.
\end{itemize}

Assembling these into a theorem, with the quantifiers in order and the finite region described
in general, is left as future work. Three worked slopes do not establish it.

\section*{Acknowledgements and AI-assistance disclosure}
\addcontentsline{toc}{section}{Acknowledgements}

Anthropic's Claude was used throughout as an assistant: for exact-arithmetic computational
design and execution, for drafting and revising this note and the certificate code, and as an
adversarial reviewer of both. Several of the mathematical repairs recorded in
\texttt{docs/CHANGELOG.md} originated in that review process, including the sign-separated
cutoff argument of Theorem~\ref{thm:cutoff}, the pairwise height certificate of
Proposition~\ref{prop:height}, and the correction of \eqref{eq:heightper}. All conceptual
choices, and final responsibility for every claim here, belong to the author. No human referee
has reviewed this note, and no part of it is formally verified.

\begin{thebibliography}{9}

\bibitem{ChaKim} B.~Cha and D.~H.~Kim.
\emph{Intrinsic Diophantine approximation on the unit circle and its Lagrange spectrum.}
Ann. Inst. Fourier (Grenoble) \textbf{73} (2023), no.~1, 101--161.
\href{https://doi.org/10.5802/aif.3522}{doi:10.5802/aif.3522}.

\bibitem{Conrad} \emph{Generalized Pell equation.} Math 154 course handout, Stanford
University. Theorem 1.1: finiteness of the equivalence classes of solutions of $x^2-dy^2=n$
under the norm-one units. Theorem 2.1: an explicit bounded region of $\mathbb{R}^2$ containing
a representative of every class.
\url{https://virtualmath1.stanford.edu/~conrad/154Page/handouts/genpell.pdf}.
The handout carries no printed author line; it is cited here by title, course and URL.

\bibitem{DeJesus2026} E.~De Jes\'us.
\emph{Boundary and Interior Record Chronologies for the Two-Mode Oscillator Spectrum.}
Zenodo (2026). \href{https://doi.org/10.5281/zenodo.21796640}{doi:10.5281/zenodo.21796640}.
Cited at repository commit \texttt{10ea1c62c6866990ded0a85a531883841d38d034}.

\bibitem{Elsner} C.~Elsner.
\emph{On Diophantine approximations with rationals restricted by arithmetical conditions.}
Fibonacci Quart. \textbf{38} (2000), no.~1, 24--33.
\url{https://www.fq.math.ca/Scanned/38-1/elsner.pdf}.

\bibitem{KLL2022} D.~H.~Kim, S.~B.~Lee and L.~Liao.
\emph{Odd-odd continued fraction algorithm.}
Monatsh. Math. \textbf{198} (2022), 323--344.
\href{https://doi.org/10.1007/s00605-022-01704-2}{doi:10.1007/s00605-022-01704-2}.
Read in the author version arXiv:2002.06310v2.

\bibitem{KLL2025} D.~H.~Kim, S.~B.~Lee and L.~Liao.
\emph{Diophantine approximation by rational numbers of certain parity types.}
Mathematika \textbf{71} (2025), e12285.
\href{https://doi.org/10.1112/mtk.12285}{doi:10.1112/mtk.12285}.
Read in the author version arXiv:2403.12341v1.

\end{thebibliography}

\appendix
\section{Machine-generated orbit tables}\label{app:tables}

The tables below are generated by \texttt{certificate/tables.py} and are reproduced verbatim;
the same data is available as CSV in \texttt{tables/}.

\textbf{Orbit identity.} Each orbit carries a stable identifier --- \texttt{P}, \texttt{T} or
\texttt{S} for $\varphi$, $1+\sqrt3$ or $\sqrt3$, followed by its index in the canonical
enumeration, which is sorted by the representative $\eta$. The identifier is necessary, because
$N_{\mathrm{rep}}$ is \emph{not} unique: at $1+\sqrt3$ thirty values each label two distinct
orbits and $N_{\mathrm{rep}}=-143$ labels four, and at $\sqrt3$ two values label two each.
Here $N_{\mathrm{rep}}$ is the defect of the \emph{normalised representative} (reduced into
$[1,\pr)$); by \eqref{eq:transport} the defect at period $p$ is $\delta^pN_{\mathrm{rep}}$, so
at $\varphi$ the sign alternates. ``first positive state'' is the first $(m,k)$ with both
coordinates positive.

\textbf{What an exclusion row asserts.} A row ``$B$ excluded by $A$ at offset $d$ from period
$P$'' asserts that for \emph{every} $p\ge P$, the state of $A$ at period $p-d$ lies below the
state of $B$ at period $p$ in height and has strictly smaller error. Accordingly $P$ is the
maximum of the error threshold \eqref{eq:thr} in $B$-period coordinates, the ordering threshold
of Proposition~\ref{prop:height} converted by $+d$, and the first period at which both states
are positive. The two states at $p=P$ are printed, so a reader can re-derive any row from
Propositions~\ref{prop:err}--\ref{prop:height} without re-running a search. Every row is
re-verified exactly at $p=P,P+1,P+2$ before it is emitted, and check C11 of
\texttt{verify\_records.py} re-derives all $67$ rows independently from the orbit each one
names and tests them at $p=P,\dots,P+11$.

\bigskip
\noindent\textbf{Classification of all admissible primitive orbits below the cutoff.}
{\small% machine-generated by certificate/tables.py -- do not edit
\begin{longtable}{llrrlll}
\hline
id & $\alpha$ & $|N|$ & $N_{\mathrm{rep}}$ & $\eta$ (rep) & first positive state & class \\
\hline
\endfirsthead
\hline
id & $\alpha$ & $|N|$ & $N_{\mathrm{rep}}$ & $\eta$ (rep) & first positive state & class \\
\hline
\endhead
\hline
\multicolumn{7}{r}{\footnotesize continued}\\
\endfoot
\hline
\endlastfoot
P0 & $\varphi$ & 1 & -1 & 1.618034 & (1,1) & strict \\
P1 & $\varphi$ & 11 & +11 & 2.381966 & (5,7) & excluded \\
P2 & $\varphi$ & 11 & -11 & 2.854102 & (3,1) & excluded \\
P3 & $\varphi$ & 5 & +5 & 3.618034 & (1,3) & alternating \\
T0 & $1+\sqrt3$ & 143 & -143 & 1.014083 & (7,9) & excluded \\
T1 & $1+\sqrt3$ & 39 & -39 & 1.052559 & (5,11) & excluded \\
T2 & $1+\sqrt3$ & 73 & +73 & 1.091035 & (3,13) & excluded \\
T3 & $1+\sqrt3$ & 193 & +193 & 1.129510 & (1,15) & excluded \\
T4 & $1+\sqrt3$ & 131 & -131 & 1.157677 & (7,11) & excluded \\
T5 & $1+\sqrt3$ & 11 & -11 & 1.196152 & (5,13) & strict \\
T6 & $1+\sqrt3$ & 111 & -111 & 1.301270 & (7,13) & excluded \\
T7 & $1+\sqrt3$ & 169 & +169 & 1.378222 & (3,17) & excluded \\
T8 & $1+\sqrt3$ & 83 & -83 & 1.444864 & (7,15) & excluded \\
T9 & $1+\sqrt3$ & 69 & +69 & 1.483340 & (5,17) & excluded \\
T10 & $1+\sqrt3$ & 47 & -47 & 1.588457 & (7,17) & excluded \\
T11 & $1+\sqrt3$ & 121 & +121 & 1.626933 & (5,19) & excluded \\
T12 & $1+\sqrt3$ & 179 & -179 & 1.693575 & (9,17) & excluded \\
T13 & $1+\sqrt3$ & 3 & -3 & 1.732051 & (1,1) & strict \\
T14 & $1+\sqrt3$ & 181 & +181 & 1.770527 & (5,21) & excluded \\
T15 & $1+\sqrt3$ & 143 & -143 & 1.837169 & (9,19) & excluded \\
T16 & $1+\sqrt3$ & 109 & +109 & 2.019238 & (7,23) & excluded \\
T17 & $1+\sqrt3$ & 47 & -47 & 2.124356 & (9,23) & excluded \\
T18 & $1+\sqrt3$ & 177 & +177 & 2.162831 & (7,25) & excluded \\
T19 & $1+\sqrt3$ & 13 & +13 & 2.267949 & (9,25) & excluded \\
T20 & $1+\sqrt3$ & 167 & -167 & 2.373067 & (11,25) & excluded \\
T21 & $1+\sqrt3$ & 107 & -107 & 2.516660 & (11,27) & excluded \\
T22 & $1+\sqrt3$ & 157 & +157 & 2.555136 & (9,29) & excluded \\
T23 & $1+\sqrt3$ & 39 & -39 & 2.660254 & (11,29) & excluded \\
T24 & $1+\sqrt3$ & 37 & +37 & 2.803848 & (11,31) & excluded \\
T25 & $1+\sqrt3$ & 183 & -183 & 2.908965 & (13,31) & excluded \\
T26 & $1+\sqrt3$ & 107 & -107 & 3.052559 & (13,33) & excluded \\
T27 & $1+\sqrt3$ & 23 & -23 & 3.196152 & (3,1) & excluded \\
T28 & $1+\sqrt3$ & 69 & +69 & 3.339746 & (13,37) & excluded \\
T29 & $1+\sqrt3$ & 191 & -191 & 3.444864 & (15,37) & excluded \\
T30 & $1+\sqrt3$ & 1 & +1 & 3.732051 & (1,3) & strict \\
T31 & $1+\sqrt3$ & 109 & +109 & 3.875644 & (15,43) & excluded \\
T32 & $1+\sqrt3$ & 191 & -191 & 3.980762 & (17,43) & excluded \\
T33 & $1+\sqrt3$ & 83 & -83 & 4.124356 & (17,45) & excluded \\
T34 & $1+\sqrt3$ & 33 & +33 & 4.267949 & (17,47) & excluded \\
T35 & $1+\sqrt3$ & 157 & +157 & 4.411543 & (17,49) & excluded \\
T36 & $1+\sqrt3$ & 183 & -183 & 4.516660 & (19,49) & excluded \\
T37 & $1+\sqrt3$ & 59 & -59 & 4.660254 & (5,1) & excluded \\
T38 & $1+\sqrt3$ & 73 & +73 & 4.803848 & (19,53) & excluded \\
T39 & $1+\sqrt3$ & 167 & -167 & 5.052559 & (21,55) & excluded \\
T40 & $1+\sqrt3$ & 121 & +121 & 5.339746 & (21,59) & excluded \\
T41 & $1+\sqrt3$ & 143 & -143 & 5.588457 & (23,61) & excluded \\
T42 & $1+\sqrt3$ & 13 & +13 & 5.732051 & (1,5) & excluded \\
T43 & $1+\sqrt3$ & 177 & +177 & 5.875644 & (23,65) & excluded \\
T44 & $1+\sqrt3$ & 111 & -111 & 6.124356 & (7,1) & excluded \\
T45 & $1+\sqrt3$ & 61 & +61 & 6.267949 & (25,69) & excluded \\
T46 & $1+\sqrt3$ & 71 & -71 & 6.660254 & (5,3) & excluded \\
T47 & $1+\sqrt3$ & 23 & -23 & 7.196152 & (3,5) & excluded \\
T48 & $1+\sqrt3$ & 181 & +181 & 7.339746 & (29,81) & excluded \\
T49 & $1+\sqrt3$ & 179 & -179 & 7.588457 & (9,1) & excluded \\
T50 & $1+\sqrt3$ & 33 & +33 & 7.732051 & (1,7) & excluded \\
T51 & $1+\sqrt3$ & 131 & -131 & 8.124356 & (7,3) & excluded \\
T52 & $1+\sqrt3$ & 97 & +97 & 8.267949 & (33,91) & excluded \\
T53 & $1+\sqrt3$ & 169 & +169 & 8.803848 & (35,97) & excluded \\
T54 & $1+\sqrt3$ & 11 & -11 & 9.196152 & (3,7) & excluded \\
T55 & $1+\sqrt3$ & 61 & +61 & 9.732051 & (1,9) & excluded \\
T56 & $1+\sqrt3$ & 143 & -143 & 10.124356 & (7,5) & excluded \\
T57 & $1+\sqrt3$ & 141 & +141 & 10.267949 & (41,113) & excluded \\
T58 & $1+\sqrt3$ & 71 & -71 & 10.660254 & (5,7) & excluded \\
T59 & $1+\sqrt3$ & 97 & +97 & 11.732051 & (1,11) & excluded \\
T60 & $1+\sqrt3$ & 193 & +193 & 12.267949 & (49,135) & excluded \\
T61 & $1+\sqrt3$ & 59 & -59 & 12.660254 & (5,9) & excluded \\
T62 & $1+\sqrt3$ & 37 & +37 & 13.196152 & (3,11) & excluded \\
T63 & $1+\sqrt3$ & 141 & +141 & 13.732051 & (1,13) & excluded \\
S0 & $\sqrt3$ & 26 & -26 & 1.124356 & (5,7) & excluded \\
S1 & $\sqrt3$ & 6 & +6 & 1.267949 & (1,3) & equality \\
S2 & $\sqrt3$ & 26 & -26 & 1.660254 & (3,1) & excluded \\
S3 & $\sqrt3$ & 22 & +22 & 1.803848 & (1,5) & excluded \\
S4 & $\sqrt3$ & 2 & -2 & 2.732051 & (1,1) & strict \\
S5 & $\sqrt3$ & 22 & +22 & 3.267949 & (3,7) & excluded \\
\end{longtable}
}

\bigskip
\noindent\textbf{Exclusion witnesses.}
{\small% machine-generated by certificate/tables.py -- do not edit
\begin{longtable}{llrlrrrlll}
\hline
id & $\alpha$ & $N_{\mathrm{rep}}$ & by & $N_A$ & $d$ & $P$ & $A$ at $P-d$ & $B$ at $P$ & rule \\
\hline
\endfirsthead
\hline
id & $\alpha$ & $N_{\mathrm{rep}}$ & by & $N_A$ & $d$ & $P$ & $A$ at $P-d$ & $B$ at $P$ & rule \\
\hline
\endhead
\hline
\multicolumn{10}{r}{\footnotesize continued}\\
\endfoot
\hline
\endlastfoot
P1 & $\varphi$ & +11 & P0 & -1 & 0 & 1 & (3,5) & (5,7) & $\Gamma_d<0$ \\
P2 & $\varphi$ & -11 & P0 & -1 & 0 & 0 & (1,1) & (3,1) & $\Gamma_d<0$ \\
T0 & $1+\sqrt3$ & -143 & T30 & +1 & 1 & 1 & (1,3) & (7,9) & $\Gamma_d<0$ \\
T1 & $1+\sqrt3$ & -39 & T30 & +1 & 1 & 1 & (1,3) & (5,11) & $\Gamma_d<0$ \\
T2 & $1+\sqrt3$ & +73 & T30 & +1 & 1 & 1 & (1,3) & (3,13) & $\Gamma_d<0$ \\
T3 & $1+\sqrt3$ & +193 & T30 & +1 & 1 & 1 & (1,3) & (1,15) & $\Gamma_d<0$ \\
T4 & $1+\sqrt3$ & -131 & T30 & +1 & 1 & 1 & (1,3) & (7,11) & $\Gamma_d<0$ \\
T6 & $1+\sqrt3$ & -111 & T1 & -39 & 0 & 1 & (5,11) & (7,13) & $\Gamma_d<0$ \\
T7 & $1+\sqrt3$ & +169 & T1 & -39 & 0 & 1 & (5,11) & (3,17) & $\Gamma_d<0$ \\
T8 & $1+\sqrt3$ & -83 & T5 & -11 & 0 & 1 & (5,13) & (7,15) & $\Gamma_d<0$ \\
T9 & $1+\sqrt3$ & +69 & T5 & -11 & 0 & 1 & (5,13) & (5,17) & $\Gamma_d<0$ \\
T10 & $1+\sqrt3$ & -47 & T5 & -11 & 0 & 1 & (5,13) & (7,17) & $\Gamma_d<0$ \\
T11 & $1+\sqrt3$ & +121 & T1 & -39 & 0 & 1 & (5,11) & (5,19) & $\Gamma_d<0$ \\
T12 & $1+\sqrt3$ & -179 & T1 & -39 & 0 & 1 & (5,11) & (9,17) & $\Gamma_d<0$ \\
T14 & $1+\sqrt3$ & +181 & T1 & -39 & 0 & 1 & (5,11) & (5,21) & $\Gamma_d<0$ \\
T15 & $1+\sqrt3$ & -143 & T1 & -39 & 0 & 1 & (5,11) & (9,19) & $\Gamma_d<0$ \\
T16 & $1+\sqrt3$ & +109 & T5 & -11 & 0 & 1 & (5,13) & (7,23) & $\Gamma_d<0$ \\
T17 & $1+\sqrt3$ & -47 & T5 & -11 & 0 & 1 & (5,13) & (9,23) & $\Gamma_d<0$ \\
T18 & $1+\sqrt3$ & +177 & T5 & -11 & 0 & 1 & (5,13) & (7,25) & $\Gamma_d<0$ \\
T19 & $1+\sqrt3$ & +13 & T13 & -3 & 0 & 1 & (7,19) & (9,25) & $\Gamma_d<0$ \\
T20 & $1+\sqrt3$ & -167 & T5 & -11 & 0 & 1 & (5,13) & (11,25) & $\Gamma_d<0$ \\
T21 & $1+\sqrt3$ & -107 & T5 & -11 & 0 & 1 & (5,13) & (11,27) & $\Gamma_d<0$ \\
T22 & $1+\sqrt3$ & +157 & T5 & -11 & 0 & 1 & (5,13) & (9,29) & $\Gamma_d<0$ \\
T23 & $1+\sqrt3$ & -39 & T13 & -3 & 0 & 1 & (7,19) & (11,29) & $\Gamma_d<0$ \\
T24 & $1+\sqrt3$ & +37 & T13 & -3 & 0 & 1 & (7,19) & (11,31) & $\Gamma_d<0$ \\
T25 & $1+\sqrt3$ & -183 & T5 & -11 & 0 & 1 & (5,13) & (13,31) & $\Gamma_d<0$ \\
T26 & $1+\sqrt3$ & -107 & T5 & -11 & 0 & 1 & (5,13) & (13,33) & $\Gamma_d<0$ \\
T27 & $1+\sqrt3$ & -23 & T13 & -3 & 0 & 1 & (7,19) & (13,35) & $\Gamma_d<0$ \\
T28 & $1+\sqrt3$ & +69 & T13 & -3 & 0 & 1 & (7,19) & (13,37) & $\Gamma_d<0$ \\
T29 & $1+\sqrt3$ & -191 & T5 & -11 & 0 & 1 & (5,13) & (15,37) & $\Gamma_d<0$ \\
T31 & $1+\sqrt3$ & +109 & T13 & -3 & 0 & 1 & (7,19) & (15,43) & $\Gamma_d<0$ \\
T32 & $1+\sqrt3$ & -191 & T5 & -11 & 0 & 1 & (5,13) & (17,43) & $\Gamma_d<0$ \\
T33 & $1+\sqrt3$ & -83 & T13 & -3 & 0 & 1 & (7,19) & (17,45) & $\Gamma_d<0$ \\
T34 & $1+\sqrt3$ & +33 & T13 & -3 & 0 & 1 & (7,19) & (17,47) & $\Gamma_d<0$ \\
T35 & $1+\sqrt3$ & +157 & T5 & -11 & 0 & 1 & (5,13) & (17,49) & $\Gamma_d<0$ \\
T36 & $1+\sqrt3$ & -183 & T5 & -11 & 0 & 1 & (5,13) & (19,49) & $\Gamma_d<0$ \\
T37 & $1+\sqrt3$ & -59 & T13 & -3 & 0 & 1 & (7,19) & (19,51) & $\Gamma_d<0$ \\
T38 & $1+\sqrt3$ & +73 & T13 & -3 & 0 & 1 & (7,19) & (19,53) & $\Gamma_d<0$ \\
T39 & $1+\sqrt3$ & -167 & T13 & -3 & 0 & 1 & (7,19) & (21,55) & $\Gamma_d<0$ \\
T40 & $1+\sqrt3$ & +121 & T13 & -3 & 0 & 1 & (7,19) & (21,59) & $\Gamma_d<0$ \\
T41 & $1+\sqrt3$ & -143 & T13 & -3 & 0 & 1 & (7,19) & (23,61) & $\Gamma_d<0$ \\
T42 & $1+\sqrt3$ & +13 & T30 & +1 & 0 & 0 & (1,3) & (1,5) & $\Gamma_d<0$ \\
T43 & $1+\sqrt3$ & +177 & T13 & -3 & 0 & 1 & (7,19) & (23,65) & $\Gamma_d<0$ \\
T44 & $1+\sqrt3$ & -111 & T13 & -3 & 0 & 1 & (7,19) & (25,67) & $\Gamma_d<0$ \\
T45 & $1+\sqrt3$ & +61 & T13 & -3 & 0 & 1 & (7,19) & (25,69) & $\Gamma_d<0$ \\
T46 & $1+\sqrt3$ & -71 & T13 & -3 & 0 & 1 & (7,19) & (27,73) & $\Gamma_d<0$ \\
T47 & $1+\sqrt3$ & -23 & T30 & +1 & 0 & 0 & (1,3) & (3,5) & $\Gamma_d<0$ \\
T48 & $1+\sqrt3$ & +181 & T13 & -3 & 0 & 1 & (7,19) & (29,81) & $\Gamma_d<0$ \\
T49 & $1+\sqrt3$ & -179 & T13 & -3 & 0 & 1 & (7,19) & (31,83) & $\Gamma_d<0$ \\
T50 & $1+\sqrt3$ & +33 & T30 & +1 & 0 & 0 & (1,3) & (1,7) & $\Gamma_d<0$ \\
T51 & $1+\sqrt3$ & -131 & T13 & -3 & 0 & 1 & (7,19) & (33,89) & $\Gamma_d<0$ \\
T52 & $1+\sqrt3$ & +97 & T13 & -3 & 0 & 1 & (7,19) & (33,91) & $\Gamma_d<0$ \\
T53 & $1+\sqrt3$ & +169 & T13 & -3 & 0 & 1 & (7,19) & (35,97) & $\Gamma_d<0$ \\
T54 & $1+\sqrt3$ & -11 & T30 & +1 & 0 & 0 & (1,3) & (3,7) & $\Gamma_d<0$ \\
T55 & $1+\sqrt3$ & +61 & T30 & +1 & 0 & 0 & (1,3) & (1,9) & $\Gamma_d<0$ \\
T56 & $1+\sqrt3$ & -143 & T30 & +1 & 0 & 0 & (1,3) & (7,5) & $\Gamma_d<0$ \\
T57 & $1+\sqrt3$ & +141 & T13 & -3 & 0 & 1 & (7,19) & (41,113) & $\Gamma_d<0$ \\
T58 & $1+\sqrt3$ & -71 & T30 & +1 & 0 & 0 & (1,3) & (5,7) & $\Gamma_d<0$ \\
T59 & $1+\sqrt3$ & +97 & T30 & +1 & 0 & 0 & (1,3) & (1,11) & $\Gamma_d<0$ \\
T60 & $1+\sqrt3$ & +193 & T13 & -3 & 0 & 1 & (7,19) & (49,135) & $\Gamma_d<0$ \\
T61 & $1+\sqrt3$ & -59 & T30 & +1 & 0 & 0 & (1,3) & (5,9) & $\Gamma_d<0$ \\
T62 & $1+\sqrt3$ & +37 & T30 & +1 & 0 & 0 & (1,3) & (3,11) & $\Gamma_d<0$ \\
T63 & $1+\sqrt3$ & +141 & T30 & +1 & 0 & 0 & (1,3) & (1,13) & $\Gamma_d<0$ \\
S0 & $\sqrt3$ & -26 & S4 & -2 & 1 & 2 & (3,5) & (5,7) & $\Gamma_d<0$ \\
S2 & $\sqrt3$ & -26 & S1 & +6 & 0 & 1 & (1,3) & (3,1) & $\Gamma_d<0$ \\
S3 & $\sqrt3$ & +22 & S1 & +6 & 0 & 1 & (1,3) & (1,5) & $\Gamma_d<0$ \\
S5 & $\sqrt3$ & +22 & S4 & -2 & 0 & 1 & (3,5) & (3,7) & $\Gamma_d<0$ \\
\end{longtable}
}

\end{document}
